%!TEX program = pdflatex
% Mathematics article -- amsart, the neutral class of the subject. Drafting
% here and converting on acceptance costs nothing, because a journal class
% changes the front matter and the page layout and leaves the body alone.
\documentclass[11pt,reqno]{amsart}

% amsart's own page is set for the AMS printed page and reads narrow on
% screen; a journal class resets it at submission.
\usepackage[margin=1in]{geometry}

% amsart already loads amsmath, amsthm and amsfonts.
\usepackage{amssymb,mathtools}
\usepackage{graphicx}
\usepackage{bm}
\usepackage{xcolor}
\usepackage{hyperref}
% Links are coloured rather than boxed: the default frames print as
% rectangles around every citation and cross-reference.
\hypersetup{colorlinks=true,
            linkcolor=blue!55!black,
            citecolor=blue!55!black,
            urlcolor=blue!55!black}
\usepackage[capitalize,nameinlink]{cleveref}

% One shared counter, so a Lemma 2.3 and a Theorem 2.4 never collide.
\theoremstyle{plain}
\newtheorem{theorem}{Theorem}[section]
\newtheorem{lemma}[theorem]{Lemma}
\newtheorem{proposition}[theorem]{Proposition}
\newtheorem{corollary}[theorem]{Corollary}
\newtheorem{conjecture}[theorem]{Conjecture}

\theoremstyle{definition}
\newtheorem{definition}[theorem]{Definition}
\newtheorem{example}[theorem]{Example}

\theoremstyle{remark}
\newtheorem{remark}[theorem]{Remark}

% amsart numbers equations straight through; section numbering matches the
% theorem counter and survives inserting a section.
\numberwithin{equation}{section}

% A few lines -- the affiliation above all -- hold unbreakable runs slightly
% wider than the measure; let TeX stretch them rather than push type into the
% margin.
\emergencystretch=2em

\title[Counterexample to phase-resonant localization]{A period-convention counterexample to phase-resonant\\localization}

\author{Zhiyuan Yao}
\address{Lanzhou Center for Theoretical Physics, Key Laboratory of Theoretical
  Physics of Gansu Province, Key Laboratory of Quantum Theory and Applications
  of MoE, Gansu Provincial Research Center for Basic Disciplines of Quantum
  Physics, Lanzhou University, Lanzhou 730000, China}
\email{yaozy@lzu.edu.cn}

\date{\today}

% AMS journals require a 2020 Mathematics Subject Classification, primary
% first; keywords are optional and printed with it.
% \subjclass[2020]{Primary 00A00; Secondary 00B00}
% \keywords{KEYWORDS}

\begin{document}

% amsart typesets the abstract from the title block, so it is
% declared inside the document and before that block.
\begin{abstract}
A conjecture on phase-resonant quasiperiodic Schr\"odinger operators asserts
that, for Lebesgue-almost every frequency and every phase, the spectral
restriction to the energies at which the Lyapunov exponent exceeds the phase
exponent $\delta(\alpha,\theta)$ is pure point with exponentially decaying
eigenfunctions, and it asserts this for every nonconstant even real-analytic
potential on $\mathbb T=\mathbb R/\mathbb Z$. Nothing in the statement
requires $1$ to be the shortest period of the potential, and we show that it
is false as a consequence. The phase exponent is a purely arithmetic
quantity, depending on the frequency and phase alone and not on the
potential, so it cannot adapt when the potential has a smaller period. Taking
$V(t)=2e^{3/2}\cos(4\pi t)$, of primitive period $1/2$, we construct for every
standard Diophantine $\alpha$ a phase whose exponent in the period-one
coordinate is $1$ while the same operator in its primitive almost-Mathieu
coordinate has exponent $2$. The Lyapunov exponent is at least $3/2$, so the stated
inequality holds at every spectral energy; but $1<e^{3/2}<e^{2}$ places the
operator on the singular-continuous side of the sharp almost Mathieu
transition, so its spectrum is purely singular continuous and has no pure
point part at all. Defining the phase exponent through the minimal
fundamental period removes the counterexample, and the corrected statement is
not addressed here.
\end{abstract}

\maketitle

\section{Introduction}
\label{sec:intro}

Let $\mathbb T=\mathbb R/\mathbb Z$, let $V:\mathbb T\to\mathbb R$ be
nonconstant, even and real-analytic, and consider the one-frequency
quasiperiodic Schr\"odinger operator $H_{V,\alpha,\theta}$ with Lyapunov
exponent $L_V(E)$ and phase exponent
\begin{equation}
  \label{eq:delta}
  \delta(\alpha,\theta)
  =\limsup_{|k|\to\infty}
   -\frac{\log\lVert2\theta+k\alpha\rVert_{\mathbb T}}{|k|}.
\end{equation}
The following has been conjectured for this class, as the generalization to
arbitrary even analytic $V$ of a theorem proved for the almost Mathieu
operator \cite{Jitomirskaya_2024ur,Jitomirskaya_2023od}:

\begin{quote}
  (S)\quad for Lebesgue-almost every $\alpha$ and every $\theta$, the
  spectral restriction of $H_{V,\alpha,\theta}$ to
  $\{E:L_V(E)>\delta(\alpha,\theta)\}$ is pure point with exponentially
  decaying eigenfunctions.
\end{quote}

This paper shows that (S) is false as written, and that the reason is a
hypothesis it does not impose.

\paragraph{The defect.} Observe what \eqref{eq:delta} does and does not
depend on. It is a statement about how well $-2\theta$ is approximated by the
orbit $k\alpha$ in $\mathbb R/\mathbb Z$: an arithmetic quantity built from
$\alpha$ and $\theta$ alone, in which the potential never appears. Statement
(S) pairs it with $L_V(E)$, which does depend on $V$. That pairing is exactly
right when $1$ is the minimal period of $V$, which is the case for the almost
Mathieu potential the statement generalizes. It is not right otherwise,
because $\mathbb R/\mathbb Z$ admits perfectly good nonconstant even
real-analytic functions of smaller period, and for those the operator is the
same operator written in the wrong coordinate.

Take
\begin{equation}
  \label{eq:V}
  V(t)=2\lambda\cos(4\pi t),
\end{equation}
which is nonconstant, even and real-analytic on $\mathbb T$ and so satisfies
every hypothesis of (S), and whose primitive period is $1/2$. Then
$V(\theta+n\alpha)=2\lambda\cos\bigl(2\pi(2\theta+2n\alpha)\bigr)$, so
$H_{V,\alpha,\theta}$ \emph{is} an almost Mathieu operator, with coupling
$\lambda$, frequency $2\alpha$ and phase $2\theta$. Its phase resonance is
governed by the arithmetic of that primitive pair, while \eqref{eq:delta}
measures the arithmetic of $(\alpha,\theta)$. The two need not agree, and
\Cref{sec:construction} constructs, for every standard Diophantine $\alpha$, a
phase at which they differ by exactly a factor of two: the stated exponent is
$1$, the primitive one is $2$.

\paragraph{The consequence.} Fix $\lambda=e^{3/2}$ in \eqref{eq:V}. The
Lyapunov exponent is then $3/2$, so $L_V(E)>1=\delta(\alpha,\theta)$ at every
spectral energy and the hypothesis of (S) holds throughout the spectrum. But
$1<e^{3/2}<e^{2}$ places the operator on the singular-continuous side of the
sharp almost Mathieu transition in its primitive coordinate, so the spectrum
is purely singular continuous and nothing in it is pure point. Since the
standard Diophantine frequencies have full measure, this negates the
almost-every-frequency, every-phase assertion of (S).

\paragraph{What this is and is not.} It is not a claim that the conjecture is
wrong in substance. It is stated as the generalization of a theorem that is
true for the almost Mathieu operator, and the counterexample exploits the one
case the generalization did not intend: a potential whose stated period is not
its shortest. Requiring $1$ to be the minimal period of $V$ --- or,
equivalently, defining the phase exponent through the minimal fundamental
period --- removes the counterexample at a stroke. What the paper establishes
is that the hypothesis is necessary and that it is currently missing, not that
the intended theorem fails. Proving the corrected statement is a separate
problem and is not attempted here.

A companion statement in the same family, asserting exact two-sided
eigenfunction tail asymptotics in the same regime, fails to the same
construction at a larger coupling; that is the subject of a separate note and
nothing below depends on it.

The setting is the global theory of one-frequency quasiperiodic operators
\cite{Avila_2015gt,Avila_2017sp,Avila_2009tt,Avila_2009al,Avila_2011ak},
resting on the dichotomy between reducibility and nonuniform hyperbolicity
\cite{Avila_2006ro}, whose
arithmetic transitions descend from the Aubry--Andr\'e duality and the early
work of Herman, Sarnak and others
\cite{AubryAndre_1980ab,Herman_1983um,Sarnak_1982sb,Eliasson_1992fs}. For
localization in this class see
\cite{Jitomirskaya_1999mi,Jitomirskaya_1994al,Jitomirskaya_1998al,Bourgain_2000on,Bourgain_2002co,Bourgain_Book,Jitomirskaya_2016lr,Avila_2008ta,Liu_2015al,Ge_2023mj}.
The opposite mechanism, in which fast approximation by periodic potentials
rules out eigenvalues altogether, goes back to Gordon \cite{Gordon_1976ps} and
produces the singular continuous half of the almost Mathieu picture
\cite{Gordon_1997da,Jitomirskaya_1994ow} --- the half the construction below
lands in. See \cite{Damanik_2016so,Simon_2000so} for surveys.

\Cref{sec:convention} fixes the notation and shows that one operator is
described by two arithmetic pairs. \Cref{sec:construction} builds the phase at
which the two disagree, and \Cref{sec:sc} derives the failure of (S).
\Cref{sec:corrected} records the hypothesis that removes it.

\section{The period convention and the two phase exponents}
\label{sec:convention}

Throughout, $\mathbb T=\mathbb R/\mathbb Z$ and
$\lVert y\rVert_{\mathbb T}=\inf_{j\in\mathbb Z}|y-j|$. A frequency
$\alpha$ is \emph{standard Diophantine} when there are $\gamma>0$ and
$\tau\ge1$ with
\begin{equation}
  \label{eq:dioph}
  \lVert h\alpha\rVert_{\mathbb T}\ \ge\ \gamma\,|h|^{-\tau}
  \qquad (h\in\mathbb Z\setminus\{0\}).
\end{equation}
For $\tau>1$ the set of $\alpha$ satisfying \eqref{eq:dioph} for some
$\gamma>0$ has full Lebesgue measure: the set of $\alpha$ with
$\lVert h\alpha\rVert_{\mathbb T}<|h|^{-\tau}$ has measure $O(|h|^{-\tau})$,
the sum over $h\ne0$ converges, and Borel--Cantelli leaves almost every
$\alpha$ in only finitely many of these sets, which are then absorbed into
$\gamma(\alpha)$.

Write $V_\lambda$ for the potential \eqref{eq:V} at coupling $\lambda>0$; it
is nonconstant, even and real-analytic on $\mathbb T$, and therefore
admissible in (S) as it is stated. Its shortest period is
$1/2$, and that single fact is the whole of the defect. Writing
\begin{equation}
  \label{eq:primitive}
  \alpha'=2\alpha \bmod 1,
  \qquad
  \theta'=2\theta \bmod 1,
\end{equation}
we have, for every lattice site $n$,
\begin{equation}
  \label{eq:same-operator}
  V_\lambda(\theta+n\alpha)
  =2\lambda\cos\bigl(4\pi(\theta+n\alpha)\bigr)
  =2\lambda\cos\bigl(2\pi(\theta'+n\alpha')\bigr),
\end{equation}
so $H_{V_\lambda,\alpha,\theta}$ is, on the same coordinates of
$\ell^2(\mathbb Z)$ and with the same eigenfunctions, the almost Mathieu
operator $H^{\mathrm{AMO}}_{\lambda,\alpha',\theta'}$ of coupling $\lambda$.

Two arithmetic pairs therefore describe one operator: the pair
$(\alpha,\theta)$ named by (S), and the primitive pair
$(\alpha',\theta')$ that governs the operator's phase resonance. Since
$\delta$ of \eqref{eq:delta} is built from its two arguments alone and never
from the potential, nothing forces the two values to agree.
\Cref{sec:construction} shows they need not: for every standard Diophantine
$\alpha$ there is a phase at which
\begin{equation*}
  \delta(\alpha,\theta)=1
  \qquad\text{while}\qquad
  \delta(\alpha',\theta')=2 .
\end{equation*}

\section{A phase with two prescribed exponents}
\label{sec:construction}

The construction below is the whole of the counterexample. Prescribing
\emph{both}
exponents simultaneously is what the argument needs; producing a single
Liouville-type phase would not separate them.

Fix a standard Diophantine $\alpha$ with constants $\gamma,\tau$ as in
\eqref{eq:dioph}. For a colour $a\in\{0,1\}$ and a signed index
$m\in\mathbb Z$ put
\begin{equation}
  \label{eq:centers}
  z_{a,m}=\frac a2-m\alpha \bmod 1,
  \qquad
  \kappa_0=1,\quad \kappa_1=2 .
\end{equation}
Colour zero carries the integer target and colour one the half-integer
target; the exponents $\kappa_a$ are the two orders to be prescribed.

\begin{lemma}[Separation of coloured centres]
  \label{lem:separation}
  Let $q$ be a denominator of a continued-fraction convergent of $\alpha$
  and let $|n-m|<q$. If the two centres $z_{a,n}$ and $z_{b,m}$ are
  distinct, then
  \begin{equation}
    \label{eq:sq}
    \lVert z_{a,n}-z_{b,m}\rVert_{\mathbb T}\ \ge\ s_q:=A\,q^{-\tau},
    \qquad
    A=\min\{\tfrac12,\ \gamma\,2^{-\tau-1}\}>0 .
  \end{equation}
\end{lemma}

\begin{proof}
  For equal colours and $0<|n-m|<q$ the best-approximation property of
  continued fractions gives
  $\lVert z_{a,n}-z_{a,m}\rVert_{\mathbb T}
   =\lVert(n-m)\alpha\rVert_{\mathbb T}>1/(2q)$.
  For different colours write $h=n-m$. If $h\ne0$ then
  \begin{equation*}
    \lVert z_{0,n}-z_{1,m}\rVert_{\mathbb T}
    =\bigl\lVert h\alpha-\tfrac12\bigr\rVert_{\mathbb T}
    \ \ge\ \tfrac12\lVert 2h\alpha\rVert_{\mathbb T}
    \ \ge\ \gamma\,2^{-\tau-1}|h|^{-\tau},
  \end{equation*}
  the middle step because $\lVert2y\rVert_{\mathbb T}\le2\lVert y-\tfrac12\rVert_{\mathbb T}$,
  and the last by \eqref{eq:dioph}. If $h=0$ the two centres are $1/2$
  apart. Since $|h|<q$, each bound is at least $A q^{-\tau}$.
\end{proof}

Fix once and for all
\begin{equation}
  \label{eq:Bc}
  B=32,\qquad c=10^{-6},
\end{equation}
and for a positive index $n$ of colour $a$ define the compact annulus
\begin{equation}
  \label{eq:annulus}
  \mathcal A_a(n)
  =\bigl\{y\in\mathbb T:\
    c\,e^{-\kappa_a n}\le\lVert y-z_{a,n}\rVert_{\mathbb T}\le e^{-\kappa_a n}
  \bigr\}.
\end{equation}
The outer radius prescribes an approximation of order $\kappa_a$ at index
$n$; the inner radius prevents the order from being exceeded, which is what
makes the resulting exponents exact rather than merely bounded below.

\begin{lemma}[Child selection]
  \label{lem:child}
  Let $Q$ be a sufficiently large convergent denominator of $\alpha$, let
  \begin{equation*}
    \Bigl\lceil\frac QB\Bigr\rceil\le\ell<\Bigl\lfloor\frac{2Q}B\Bigr\rfloor,
  \end{equation*}
  let $\mathcal A_b(\ell)$ be the current parent, and put
  $M=\lfloor Q/2\rfloor$. Then for either colour $a$ and every
  sufficiently large later convergent denominator $q$, chosen also with
  $q/B>M$, there is an index
  \begin{equation}
    \label{eq:child-index}
    \Bigl\lceil\frac qB\Bigr\rceil\le n<\Bigl\lfloor\frac{2q}B\Bigr\rfloor
  \end{equation}
  with
  \begin{equation}
    \label{eq:nested}
    \mathcal A_a(n)\subset\mathcal A_b(\ell)
  \end{equation}
  and such that every $y\in\mathcal A_a(n)$ satisfies
  \begin{equation}
    \label{eq:protected}
    \lVert y-z_{u,m}\rVert_{\mathbb T}\ \ge\ c\,e^{-\kappa_u|m|}
    \qquad
    \bigl(u\in\{0,1\},\ M\le|m|<\lfloor q/2\rfloor\bigr).
  \end{equation}
\end{lemma}

For the selected pair $(u,m)=(a,n)$ itself, \eqref{eq:protected} is the
inner inequality of \eqref{eq:annulus}, so the child's own centre is not an
obstruction.

\begin{proof}
  Write $r=e^{-\kappa_b\ell}$ for the parent's outer radius and trim the
  parent to
  \begin{equation}
    \label{eq:trimmed}
    \mathcal U=\bigl\{y:2cr<\lVert y-z_{b,\ell}\rVert_{\mathbb T}<(1-c)r\bigr\},
  \end{equation}
  a set of length $2(1-3c)r$. Irrational rotations are uniquely ergodic,
  uniformly in the initial point of the orbit block, so the block of
  consecutive candidate indices \eqref{eq:child-index} contains at least
  \begin{equation}
    \label{eq:count}
    \frac{(1-3c)\,r\,q}{B}
  \end{equation}
  centres lying in $\mathcal U$ once $q$ is large enough. Enlarging $q$
  further makes the largest child radius at most $e^{-q/B}<cr$, so a child
  centred in $\mathcal U$ satisfies \eqref{eq:nested}.

  Call a candidate \emph{bad} when its child meets one of the obstacle balls
  $B(z_{u,m},\,c\,e^{-\kappa_u|m|})$ of \eqref{eq:protected} other than its
  own inner ball. Candidate indices obey $n<2q/B$ and obstacle indices
  $|m|<q/2$, so
  \begin{equation}
    \label{eq:index-gap}
    0<|n-m|<\frac{2q}{B}+\frac q2<q
  \end{equation}
  unless the two indices coincide, and the coinciding case with opposite
  colours is covered by \Cref{lem:separation} as well. Choose $q$ so that the
  child radius is below $s_q/4$; by \Cref{lem:separation} an obstacle meeting
  the child must then have radius at least $3s_q/4$, and since obstacle radii
  decay exponentially in $|m|$ only $O(\log q)$ signed obstacles can do so.

  Candidate centres of the fixed colour are more than $1/(2q)$ apart, so a
  single obstacle of radius $\rho$ spoils at most $4q(\rho+e^{-q/B})+1$
  candidates. Summing over both colours and both signs,
  \begin{equation}
    \label{eq:bad}
    \#\{\text{bad candidates}\}
    \ \le\ C_0\,q\,e^{-M}+O\bigl(q e^{-q/B}\log q+\log q\bigr),
  \end{equation}
  with $C_0$ independent of $Q$ and $q$. Because $\ell<2Q/B$, $\kappa_b\le2$
  and $M=\lfloor Q/2\rfloor$,
  \begin{equation}
    \label{eq:ratio}
    \frac{e^{-M}}{r}
    \ \le\ \exp\Bigl(-\frac Q2+\frac{4Q}{B}+O(1)\Bigr)
    \ =\ e^{-3Q/8+O(1)} .
  \end{equation}
  Taking the initial parent scale large enough makes the first term of
  \eqref{eq:bad} smaller than half the count \eqref{eq:count}; for that fixed
  parent the remaining terms of \eqref{eq:bad} are $o(rq)$ as $q\to\infty$.
  Hence some candidate is not bad, and it is the required child.
\end{proof}

\begin{proposition}[Two prescribed exact orders]
  \label{prop:two-orders}
  Let $\alpha$ satisfy \eqref{eq:dioph}. There is $x\in\mathbb T$ with
  \begin{equation}
    \label{eq:order1}
    \limsup_{|m|\to\infty}
      \frac{-\log\lVert x+m\alpha\rVert_{\mathbb T}}{|m|}=1
  \end{equation}
  and
  \begin{equation}
    \label{eq:order2}
    \limsup_{|m|\to\infty}
      \frac{-\log\lVert x+m\alpha-\tfrac12\rVert_{\mathbb T}}{|m|}=2,
  \end{equation}
  both orders being attained along subsequences of positive indices.
\end{proposition}

\begin{proof}
  Begin with a sufficiently large annulus of colour zero and apply
  \Cref{lem:child} repeatedly, alternating the requested colours
  $0,1,0,1,\dots$ and obtaining convergent denominators $Q_1<Q_2<\cdots$.
  This produces nested nonempty compact annuli, so their intersection
  contains a point $x$; put $M_t=\lfloor Q_t/2\rfloor\to\infty$.

  The protected windows $M_t\le|m|<M_{t+1}$ of \eqref{eq:protected} cover
  every sufficiently large signed integer, so
  \begin{equation}
    \label{eq:lower}
    \lVert x-z_{a,m}\rVert_{\mathbb T}\ \ge\ c\,e^{-\kappa_a|m|}
    \qquad(a=0,1)
  \end{equation}
  for all large $|m|$, which after applying $-|m|^{-1}\log$ bounds both
  upper limits from above by $\kappa_a$, the constant $c$ disappearing in
  the limit. Membership in the infinitely many selected annuli of each
  colour gives the matching
  $\lVert x-z_{a,n_t}\rVert_{\mathbb T}\le e^{-\kappa_a n_t}$ along
  subsequences of positive indices, hence the lower bounds. With
  $z_{0,m}=-m\alpha$ and $z_{1,m}=\tfrac12-m\alpha$ these are exactly
  \eqref{eq:order1} and \eqref{eq:order2}.
\end{proof}

\begin{corollary}
  \label{cor:exponents}
  Let $\alpha$ satisfy \eqref{eq:dioph} and let $\theta$ satisfy
  $2\theta=x\bmod1$ for the $x$ of \Cref{prop:two-orders}. Then, with
  $\alpha',\theta'$ as in \eqref{eq:primitive},
  \begin{equation}
    \label{eq:two-deltas}
    \delta(\alpha,\theta)=1,
    \qquad
    \delta(\alpha',\theta')=2,
  \end{equation}
  and $\alpha'$ is again standard Diophantine.
\end{corollary}

\begin{proof}
  Since $2\theta=x$, \eqref{eq:delta} and \eqref{eq:order1} give
  $\delta(\alpha,\theta)=1$. For the primitive pair use the exact circle
  identity
  \begin{equation}
    \label{eq:doubling}
    \lVert 2y\rVert_{\mathbb T}
    =2\min\bigl\{\lVert y\rVert_{\mathbb T},\
                 \lVert y-\tfrac12\rVert_{\mathbb T}\bigr\},
  \end{equation}
  applied at $y=x+k\alpha$. Because $2\theta'+k\alpha'=2(x+k\alpha)$ modulo
  $1$,
  \begin{equation*}
    \delta(\alpha',\theta')
    =\max\left\{
      \limsup_{k\to+\infty}\frac{-\log\lVert x+k\alpha\rVert_{\mathbb T}}{k},\
      \limsup_{k\to+\infty}
        \frac{-\log\lVert x+k\alpha-\tfrac12\rVert_{\mathbb T}}{k}
    \right\}=2,
  \end{equation*}
  the additive $-\log2/k$ from \eqref{eq:doubling} vanishing in the limit and
  the two upper limits being $1$ and $2$ by \Cref{prop:two-orders}. Finally
  $\lVert k\alpha'\rVert_{\mathbb T}=\lVert2k\alpha\rVert_{\mathbb T}
   \ge\gamma(2|k|)^{-\tau}$, so $\alpha'$ satisfies \eqref{eq:dioph} with
  constant $\gamma2^{-\tau}$.
\end{proof}

\begin{remark}
  \label{rem:index-set}
  The two exponents are the same under either index convention. Sources
  differ on whether the upper limit in \eqref{eq:delta} is taken over
  $k\to+\infty$ or over $|k|\to\infty$, and in a paper about a convention
  that difference cannot be left implicit. \Cref{prop:two-orders} settles it:
  the lower bounds come from subsequences of positive indices, so they hold
  for the one-sided reading, while \eqref{eq:lower} bounds the quotient above
  for every sufficiently large $|m|$ of either sign, so the two-sided upper
  limit is no larger. Both readings therefore return
  $\delta(\alpha,\theta)=1$ and $\delta(\alpha',\theta')=2$, and nothing below
  depends on the choice.
\end{remark}

\begin{proposition}[The Lyapunov exponent is carried across]
  \label{prop:lyapunov}
  For every $\lambda>0$ and every energy $E$,
  \begin{equation}
    \label{eq:lyap}
    L_{V_\lambda,\alpha}(E)=L_{\mathrm{AMO},\alpha'}(E),
  \end{equation}
  where the right-hand side is the Lyapunov exponent of the almost Mathieu
  cocycle of coupling $\lambda$ and frequency $\alpha'$.
\end{proposition}

\begin{proof}
  By \eqref{eq:same-operator} the two transfer matrices agree after doubling
  the base point: if $A^{V_\lambda,\alpha}_E(N,t)$ and
  $A^{\mathrm{AMO},\alpha'}_E(N,s)$ denote the $N$-step matrices, then
  $A^{V_\lambda,\alpha}_E(N,t)=A^{\mathrm{AMO},\alpha'}_E(N,2t)$ for every
  $t$. The doubling map $t\mapsto2t$ preserves Haar measure on $\mathbb T$,
  so the two integrals of $N^{-1}\log\lVert A_E(N,\cdot)\rVert$ coincide for
  every $N$, and so do their limits.
\end{proof}

\section{Singular continuous spectrum where localization was claimed}
\label{sec:sc}

Take $\lambda=e^{3/2}$ in \eqref{eq:V} and let $\alpha,\theta$ be as in
\Cref{cor:exponents}. We use the sharp almost Mathieu transition of
Jitomirskaya and Liu \cite{Jitomirskaya_2024ur}, whose normalization is
fixed by the transfer matrix
\begin{equation}
  \label{eq:amo-normalization}
  A(t)=
  \begin{pmatrix}
    E-2\lambda\cos(2\pi t) & -1\\
    1 & 0
  \end{pmatrix},
\end{equation}
the convention in which \eqref{eq:same-operator} was written. Their Theorem
1.1 gives purely singular continuous spectrum for the Diophantine almost
Mathieu operator when $1<|\lambda|<e^{\delta(\alpha',\theta')}$.

\begin{theorem}
  \label{thm:main}
  Statement (S) is false. For $V_{e^{3/2}}(t)=2e^{3/2}\cos(4\pi t)$ and
  every standard Diophantine $\alpha$ there is a phase $\theta$ such that
  $L_{V,\alpha}(E)>\delta(\alpha,\theta)$ at every $E$ in the spectrum,
  while the spectrum of $H_{V,\alpha,\theta}$ is purely singular continuous
  and hence has no pure point part at all.
\end{theorem}

\begin{proof}
  Fix $\theta$ as in \Cref{cor:exponents}. By \eqref{eq:same-operator} the
  operator is $H^{\mathrm{AMO}}_{e^{3/2},\alpha',\theta'}$, and by
  \Cref{cor:exponents} the pair $(\alpha',\theta')$ is Diophantine with
  $\delta(\alpha',\theta')=2$. Since
  \begin{equation}
    \label{eq:sandwich}
    1<e^{3/2}<e^{2}=e^{\delta(\alpha',\theta')},
  \end{equation}
  Theorem 1.1 of \cite{Jitomirskaya_2024ur} applies and the spectrum is
  purely singular continuous: the spectral measures carry no atoms, so no
  spectral restriction of the operator is pure point.

  For the hypothesis of (S), \Cref{prop:lyapunov} carries the question to the
  almost Mathieu cocycle of coupling $e^{3/2}$, for which Herman's
  subharmonicity bound gives $L(E)\ge\log\lambda=3/2$ at every energy
  \cite{Herman_1983um,Damanik_2016so}; the bound is in fact an equality on the
  spectrum \cite{Bourgain_2002co}, though only the inequality is used here.
  Meanwhile $\delta(\alpha,\theta)=1$ by \Cref{cor:exponents}. Hence
  \begin{equation}
    \label{eq:hyp-holds}
    L_{V,\alpha}(E)\ \ge\ \tfrac32\ >\ 1=\delta(\alpha,\theta)
  \end{equation}
  throughout the spectrum, so the restriction (S) speaks about is the whole
  of it, and (S) asserts that it is pure point with exponentially decaying
  eigenfunctions. The two conclusions are mutually exclusive.

  The standard Diophantine frequencies have full Lebesgue measure, so the
  frequencies exhibited here cannot be absorbed into the null set that the
  phrase ``for Lebesgue-almost every $\alpha$'' permits.
\end{proof}

\section{The corrected statement}
\label{sec:corrected}

\begin{remark}
  \label{rem:repairs}
  Either of the following restores the intended content of (S). Require $1$
  to be the shortest period of $V$; then $(\alpha,\theta)=(\alpha',\theta')$
  in the notation of \eqref{eq:primitive} and the two exponents coincide by
  construction. Alternatively, replace $\delta$ in \eqref{eq:delta} by a
  resonance exponent invariant under the finite covering reparametrizations
  $t\mapsto kt$ and computed over every symmetry centre of $V$, which for
  $V_\lambda$ returns the value $2$ that governs the operator.
\end{remark}

It is worth saying which half of the almost Mathieu picture needs the
normalization and which does not. Jitomirskaya and Liu also prove, for every
even Lipschitz $v$ on $\mathbb T$, every irrational $\alpha$ and every
$\theta$, that $H_{v,\alpha,\theta}$ has no eigenvalues where
$L(E)<\delta(\alpha,\theta)$, with $\delta$ given by \eqref{eq:delta} and
with no hypothesis on the period of $v$ \cite{Jitomirskaya_2024ur}. That
statement is immune to the reparametrization used here, because passing to the
primitive coordinate can only increase $\delta$ and leaves $L$ alone, so
$L<\delta$ is preserved. It is the opposite regime, $L>\delta$, that (S)
asserts and that the present counterexample occupies: there the inequality
points the other way, and an understated $\delta$ turns a true hypothesis
into a false conclusion.

The repaired statement is not addressed here, and nothing above bears on
whether it is true. What \Cref{thm:main} establishes is that the hypothesis
is necessary and that it is currently absent: the almost Mathieu theorem of
\cite{Jitomirskaya_2024ur,Jitomirskaya_2023od} that (S) generalizes is stated
for a potential whose shortest period is $1$, and the generalization to
arbitrary even analytic $V$ carried $\delta$ across without carrying that
normalization with it.

\section*{Acknowledgments}

Z.\ Y.\ acknowledges support from the Strategic Priority Research Program of the
Chinese Academy of Sciences under Grant No.~XDB1680102. Z.\ Y.\ also
acknowledges support by National Natural Science Foundation of China (Grant
No.~12247101). We also acknowledge Gewu Intelligence Lab for providing the
collaborative research environment in which this project was developed.

Large language models were used substantively in this work: OpenAI GPT-5.6 Sol
in solving the problem, including the formulation of arguments and the
derivation of intermediate steps, and Anthropic Claude Opus~5 in surveying the
literature, in verifying the period convention against the source statements,
in checking every attribution against the cited sources, and in drafting and
revising the manuscript. The author specified the problem,
guided the approach to be taken, checked every argument and result they
produced, and takes full responsibility for the content of this manuscript.

\bibliographystyle{amsplain}
\bibliography{references,unverified}

\end{document}
